% !PRINTED 420
\weberpage{420}
Ist $\tau$ eine Variable mit positiv imaginärem Teil, und $v$ gleichfalls eine Variable, so zerfällt $F_n[v, j(\tau)]$ nach § 69 in die linearen Faktoren

\[
v - j\left(\frac{c + \partial\,\tau}{a + b\,\tau}\right),
\]

worin $\left(\begin{smallmatrix}a, & b\\ c, & \partial\end{smallmatrix}\right)$ ein vollständiges Repräsentantensystem nicht äquivalenter Transformationen $n$ten Grades durchläuft, und folglich kann $F_n(u, u)$ in die Faktoren

\weq{5}{j(\tau) - j\left(\frac{c + \partial\,\tau}{a + b\,\tau}\right)}

zerlegt werden, wenn $u = j(\tau)$ gesetzt wird.

Aus dem Theorem II., § 114, aber folgt, daß, wenn wir $\tau = \omega$ setzen, $k$ von den Faktoren (5) verschwinden. Wir können aber auch noch folgern, daß der Quotient

\weq{6}{\dfrac{j(\tau) - j\left(\dfrac{c + \partial\,\tau}{a + b\,\tau}\right)}{j(\tau) - j(\omega)}}

für $j(\tau) = j(\omega)$ endlich und von Null verschieden bleibt. Denn differentiiert man Zähler und Nenner nach $\tau$ und setzt dann $\tau = \omega$, so ergibt sich der Grenzwert:

\[
1 - \frac{n}{(a + b\,\omega)^2} = \frac{2\sqrt{D}\,x}{y + \sqrt{D}\,x},
\]

der von Null verschieden ist.

Die Invariante $j(\omega)$ bleibt unverändert, wenn $\omega$ durch eine äquivalente Zahl ersetzt wird, hat aber für jede andere Zahl $\omega$ einen anderen Wert. Sie gehört also nicht zu der individuellen Form $(A, B, C)$, sondern zu der ganzen \gesperrt{Klasse} äquivalenter Formen, und wird darum die \gesperrt{Klasseninvariante} genannt.

Ferner ist die Zahl $k$ nicht von den Koeffizienten $A, B, C$, sondern nur von $D$ und $n$ abhängig, und ist daher für alle primitiven Formen der Diskriminante $D$ dieselbe. Wir führen nun die folgende Funktion der Variablen $u$ ein: Es bedeute $\omega_1, \omega_2, \ldots, \omega_h$ ein vollständiges System nicht äquivalenter Zahlen der Diskriminante $D$, also das System der Wurzeln eines Systems nicht äquivalenter Formen $(A, B, C)$ der Diskriminante $D$. Wir setzen

\weq{7}{H_m(u) = [u - j(\omega_1)]\,[u - j(\omega_2)]\,\ldots\,[u - j(\omega_h)],}
%CONT>>
